% =====================================================================================================
% Section V: Economic Implications. Writer: economics writer (v3), 2026-09-24; round-3 revision by WP5, 2026-09-25
% (fix list E1-E3: 49 budget-level decision units, trend variants, fleet share moved to Online Appendix F).
% Module rb3_econ2 (memo output/memos/rb3_econ2.md; round-3 changes in output/memos/rb5_econ.md).
% Table 3 = paper/tables/table3_econ.tex (copy of output/tables/rb3_econ2_table.tex). The economics figure and three
% supporting tables are in the Online Appendix (sections/appendix_econ.tex). Numbers and sources:
% paper/notes/round3_WP5_log.md. Version 2 of this file: sections/v2/economics.tex.
% =====================================================================================================

\section{Economic Implications}\label{sec:econ}

Three economic questions turn on the objects of Sections~\ref{sec:tech} and~\ref{sec:wedge}: how fast frontier
training's demand for data grows, what a limit on unique text costs, and how the returns to compute should enter models
of growth. The answers load on different objects, and the object on which the designs agree most closely matters least.
Compute-optimal data demand depends on the allocation exponent $a$ and the compute-optimal ratio $M^*(C)$, not on
$\sigma^*$; the cost of a data wall depends mainly on how fast repeated data lose value. The curvature sets the scale of
the revealed value of compactness, and it converts that value into data: at given compute $D/D^*(C)=[M/M^*(C)]^{1/2}$,
so by equation~\eqref{eq:wM} a model with wedge $w$ processes
\begin{equation}\label{eq:econ-DD}
\frac{D}{D^*(C)}=w^{\sigma^*/[2(1-\sigma^*)]}
\end{equation}
times the compute-optimal tokens.

\subsection{The Growth of Data Demand}\label{sec:econ-demand}

Along the expansion path, compute-optimal data grow as $D^*(C)\propto C^{1-a}$. Frontier training compute grew
5.06-fold a year from 2018 to 2026 on the current Epoch AI database (95 percent confidence interval 4.27 to 6.02),
faster than the 4.2-fold rate that \citet{sevilla2024training} published for 2018 to May 2024; the difference is a
vintage effect, and we report both rates.\footnote{We regress $\log_{10}C$ on release date for the running top ten of training compute at release in \citet{epochai2026data} (92 runs by 27 developers), with a wild
cluster bootstrap-$t$ interval by developer \citep{webb2023reworking}. On Epoch's public file of May 31, 2024
\citep{epochai2024dataarchive}, the same estimator gives 4.13-fold for January 2018 to May 2024, against 5.23-fold on
the current file for the same window; revised compute estimates explain 64 percent of that gap and entries added later
the rest (Online Appendix Table~\ref{tab:app-econ-growth}).} Across technologies, $a$ runs from 0.36 (Farseer's own
form, beyond its design) to 0.57 (the RefinedWeb sweep of \citealt{gadre2024language}), so compute-optimal data demand
grows 2.0-fold to 2.8-fold a year at 5.06-fold compute growth and 1.9-fold to 2.5-fold at the published rate
(Table~\ref{tab:econ}, panel~A). The top of these ranges rests on two imprecise slopes, Farseer's and the OLMo
ladder's; the model-free IsoFLOP paths of Section~\ref{sec:tech} give 2.3-fold to 2.9-fold. Levels are not portable
either: at the fitted frontier compute of September 2026, $9.4\times10^{26}$ FLOP, compute-optimal data range from 9
to 519 trillion tokens.

We therefore anchor each technology at the observed data use of the four disclosed dense training runs above $10^{25}$
FLOP, whose median lies close to the expansion path ($w=1.21$ under the reference). Frontier runs then reach the
quality-adjusted stock of public text, about 100 trillion tokens \citep{villalobos2022run}, between 2026.6 and 2027.8
across technologies at 5.06-fold compute growth, and the repetition-adjusted stock of 320 trillion tokens between 2027.8
and 2029.6 (2028.0 and 2030.0 at the published rate). The stock's size matters more than $a$: at its lower bound,
today's frontier runs already process 2.2 to 4.9 times the stock.

\textit{Rising wedges.} These forecasts hold the frontier's data multiple $D/D^*$ fixed, but the wedges revealed by
open-weight releases rose. Under the reference technology, a regression of $\ln w$ on release date over the 49 decision
units gives growth of 1.86-fold a year from 2023 to 2025 (95 percent confidence interval 0.96 to 3.54). If the frontier
adopts this trend, a scenario across populations since the trend comes mostly from small models,
equation~\eqref{eq:econ-DD} makes its data multiple grow 2.06-fold a year at $\sigma^*=0.70$, and frontier data demand
4.6-fold a year rather than 2.2-fold (Table~\ref{tab:econ}, panel~B). At the point estimate the wedge trend adds about
as much to the growth of data demand as compute growth does, but its interval includes no growth, as do those of the
trends within developers, above $10^{24}$ FLOP and weighted by compute (1.59-fold, 1.76-fold and 2.21-fold a year).
Extrapolating the trend in $M/M^*(C)$ instead needs no curvature and gives 1.97-fold to 2.38-fold a year for the data
multiple over the 77 models under the 32 technologies (2.24 under the reference); the $\sigma^*$ rows of panel~B
convert a projected value of compactness into data and are not a forecast interval. If the rise stops at the 2025
compute-weighted wedge of the decision units, 5.38, which rests largely on Alibaba's Qwen3, five-year frontier data
demand is 171 to 468 times today's (114 to 215 without Alibaba), against 56 times with the wedge held.

\subsection{The Data Wall}\label{sec:econ-wall}

A data wall concerns unique tokens $U$, whereas $D$ counts tokens processed, including repetitions. We combine the
technology with the repetition model of \citet{muennighoff2023scaling}, in which effective data are
$D'=U+UR^*_D(1-e^{-R/R^*_D})$ with $R=D/U-1$ repetitions: the data-decay scale $R^*_D=15.4$ alone is the most favorable
case, their data-only fit has $R^*_D=2.9$, and in their full model excess parameters also lose value. The cost of the
wall is the increase in lifetime cost, training plus a value of compactness proportional to $N$ that reproduces the
developer's wedge, at the loss of its unconstrained choice; for a compute-optimal run it depends on scarcity
$r\equiv D^*(C)/U$ alone. At $r=4$ it is
7.7, 7.4 and 7.2 percent for $\sigma^*=0.60$, 0.70 and 0.74 if repeated data decay slowly, but 51, 40 and 35 percent
under the data-only fit (Table~\ref{tab:econ}, panel~C). The repetition technology, not the curvature, drives the cost of the wall.

Over-training brings the wall forward. A model with the 2025 aggregate wedge processes about 7.2 times the
compute-optimal tokens at given compute, and at $r=4$ it pays 16 percent more lifetime cost rather than 7.4 (Online
Appendix Table~\ref{tab:app-econ-wall}). Today's frontier runs, at their observed data use, process 0.5 to 1.1 times the
100-trillion-token stock, so the wall costs them less than 0.1 percent. It binds within the forecast horizon: with the
wedge held and 5.06-fold compute growth, frontier runs process 3.5 times the stock in 2029, and the wall costs 4.6
percent of lifetime cost if repeated data decay slowly (about \$2 billion at \$$10^{-18}$ per FLOP, the median amortized
cost of 2024--2025 frontier runs; \citealp{epochai2026data}) and 23 percent under the data-only fit. With the wedge
rising, they process 18 times the stock, at a cost of 10.4 or 48 percent. At the 2026 frontier, the shadow value of a
unique token exceeds posted prices of generated tokens \citep{deepseek2026pricing} once compute-optimal data exceed 1.5
to 2.9 times the stock, a lower bound because synthetic tokens are imperfect substitutes (Online Appendix
Table~\ref{tab:app-econ-gamma}).

\textit{Feedback on measured wedges.} A binding cap on tokens adds a term to the denominator of the wedge
(Proposition~\ref{prop:wedge}(iii)), so measured wedges understate the value of compactness as the wall approaches. When
a developer with the 2025 aggregate wedge would process four times the unique stock, its measured wedge falls short of
the wedge it would reveal without the wall by 4 percent if repeated data decay slowly, by 15 percent under the data-only
fit and by 53 percent under a hard cap (Online Appendix Table~\ref{tab:app-econ-wall}, panel~B). As releases approach
the stock, measured wedges should therefore flatten even if the value of compactness keeps rising, so a flattening among
2026 and 2027 releases would not by itself show that compactness has become less valuable (except under the
full repetition model; Online Appendix~\ref{app:econ}).

% Table 3 (round 3): copied from output/tables/rb3_econ2_table.tex (module rb3_econ2), placed by WP5, 2026-09-25.
% Regenerate with code/analysis/rb3_econ2/run.py and re-copy (the copy adds only these two comment lines).
% Table 3 (revised) -- Economic implications: data demand, the data wall and the revealed wedge (Section V).
% Module rb3_econ2 (code/analysis/rb3_econ2/run.py); sources: rb3_econ2_demand.csv, _scenarios.csv,
% _wall_compute_optimal.csv, _frontier_wall.csv, _wedge_by_year.csv, _wedge_trend.csv, _fleet.csv, _gamma.csv.
\begin{table}[tp]
\centering
\caption{Economic Implications: Data Demand, the Data Wall and the Revealed Wedge}
\label{tab:econ}
\footnotesize
\setlength{\tabcolsep}{1.8pt}
\renewcommand{\arraystretch}{0.96}
\begin{tabular*}{\textwidth}{@{\extracolsep{\fill}}lccccccc@{}}
\toprule
\multicolumn{8}{@{}l}{\textit{Panel A. Compute-optimal data demand}} \\[2pt]
 & & $D^*$ in & \multicolumn{2}{c}{$D^*$ per year} & \multicolumn{2}{c}{$D^*$ over 5 years} & 100T \\
 & & 2026 & \multicolumn{2}{c}{if $C$ grows} & \multicolumn{2}{c}{if $C$ grows} & reached \\
\cmidrule(lr){4-5}\cmidrule(lr){6-7}
Technology & $a$ & (T) & $4.2\times$ & $5.06\times$ & $4.2\times$ & $5.06\times$ & (year) \\
\midrule
Chinchilla, $\kappa$ free (reference) & 0.50 & 57 & 2.04 & 2.24 & 35 & 56 & 2027.4 \\
Chinchilla, $\kappa=1$ & 0.51 & 48 & 2.01 & 2.20 & 33 & 52 & 2027.4 \\
Llama 3 (Meta's law) & 0.46 & 90 & 2.16 & 2.39 & 47 & 78 & 2027.1 \\
Farseer, own form (local) & 0.36 & 419 & 2.52 & 2.84 & 101 & 183 & 2026.6 \\
Gadre et al., RefinedWeb & 0.57 & 9.3 & 1.87 & 2.02 & 23 & 34 & 2027.8 \\
OLMo ladder (AI2) & 0.37 & 519 & 2.47 & 2.78 & 92 & 165 & 2026.7 \\
\midrule
\multicolumn{8}{@{}l}{\textit{Panel B. Scenario: the frontier adopts the open-weight trend}} \\[2pt]
 & & $D/D^*$ & \multicolumn{2}{c}{$D$ per year} & \multicolumn{2}{c}{$D$ over 5 years$^a$} & 100T \\
 & & per & \multicolumn{2}{c}{if $C$ grows} & \multicolumn{2}{c}{if $C$ grows} & reached \\
\cmidrule(lr){4-5}\cmidrule(lr){6-7}
Trend: $w$ per year [95\% CI] & $\sigma^*$ & year & $4.2\times$ & $5.06\times$ & $4.2\times$ & $5.06\times$ & (year) \\
\midrule
None: wedge held ($w=1.21$) & -- & 1.00 & 2.04 & 2.24 & 35 & 56 & 2027.4 \\
Decision units: 1.86 [0.96, 3.54] & 0.60 & 1.59 & 3.24 & 3.56 & 108 & 171 & 2027.1 \\
 & 0.70 & 2.06 & 4.20 & 4.61 & 201 & 319 & 2027.1 \\
 & 0.74 & 2.41 & 4.92 & 5.40 & 294 & 468 & 2027.0 \\
Within developers: 1.59 [0.47, 5.59] & 0.70 & 1.72 & 3.51 & 3.85 & 201 & 319 & 2027.1 \\
Units above $10^{24}$ FLOP: 1.76 [0.78, 4.15] & 0.70 & 1.93 & 3.94 & 4.32 & 201 & 319 & 2027.1 \\
Compute-weighted units: 2.21 [0.38, 13.1] & 0.70 & 2.52 & 5.14 & 5.64 & 201 & 319 & 2027.0 \\
Flagships above $10^{25}$ FLOP: 1.70$^b$ & 0.70 & 1.86 & 3.78 & 4.15 & 201 & 319 & 2027.1 \\
\midrule
\multicolumn{8}{@{}l}{\textit{Panel C. The data wall at compute-optimal and at observed allocations}} \\[2pt]
 & & & \multicolumn{3}{c}{Extra cost (\%), repetition:} & Shadow & Measured \\
 & & & $R^*_D=$ & $R^*_D=$ & Full & value & wedge \\
\cmidrule(lr){4-6}
 & $\sigma^*$ & Scarcity & 15.4 & 2.9 & model & (\$/M) & ratio$^c$ \\
\midrule
\multicolumn{8}{@{}l}{\quad\textit{Compute-optimal run ($w=1$), $10^{26}$ FLOP; scarcity $r=D^*(C)/U$}} \\
\quad $r=4$ & 0.60 & 4 & 7.7 & 51.4 & 42.0 & 3.0 & 0.93 \\
 & 0.70 & 4 & 7.4 & 39.8 & 40.3 & 2.8 & 0.92 \\
 & 0.74 & 4 & 7.2 & 35.3 & 39.5 & 2.7 & 0.91 \\
\multicolumn{8}{@{}l}{\quad\textit{Projected frontier at its observed allocation, 100T stock; scarcity $D/U$}} \\
\quad 2029, wedge held ($w=1.21$) & 0.70 & 3.5 & 4.6 & 23.0 & 13.2 & 31.4 & 0.93 \\
\quad 2029, wedge rises ($w=4.95$) & 0.70 & 17.9 & 10.4 & 48.3 & 10.4 & 175 & 0.71 \\
\bottomrule
\end{tabular*}
\vspace{2pt}
\begin{tablenotes}[Notes]
$C$ grows 4.2-fold a year (Epoch's published rate) or 5.06-fold (our estimate). Panel A: $a$, allocation exponent ($D^*\propto C^{1-a}$); Farseer's own form: local slope at $10^{22}$--$10^{23}$ FLOP, beyond its design; OLMo: bootstrap interval for $a$ [0.11, 0.82]; model-free IsoFLOP paths ($a=0.35$--0.50): 2.06--2.53 and 2.26--2.86-fold a year; $D^*$ at the fitted frontier compute of September 2026 ($9.4\times10^{26}$ FLOP); 100T reached (at $5.06\times$): year frontier runs, at the median $D/D^*$ of the four disclosed dense runs above $10^{25}$ FLOP, reach the quality-adjusted stock of public text (100T tokens in 2024, growing 5 percent a year). Panel B: from September 2026 the frontier's wedge ($w=1.21$, $D/D^*=1.25$) grows at an open-weight trend of 2023--2025 (OLS of $\ln w$ on release date, with 95 percent wild cluster bootstrap-$t$ intervals by developer; 6 developers above $10^{24}$ FLOP; Online Appendix Table~\ref{tab:app-econ-trend}) and $D/D^*=w^{\sigma^*/[2(1-\sigma^*)]}$; tokens per parameter give the $\sigma^*=0.70$ row whatever $\sigma^*$ (1.97--2.38-fold a year across 32 technologies). $^a$The rise stops at the 2025 wedge ($w=5.38$; 3.11 without Alibaba's Qwen3), reached after 1.9 to 3.2 years. $^b$Four disclosed dense runs above $10^{25}$ FLOP, 2024--2025; descriptive (1.31 without Pangu Ultra). Panel C: the developer minimizes training cost plus a value of compactness proportional to $N$ at the loss of its uncapped choice; extra cost relative to that choice; repetition model of \citet{muennighoff2023scaling} (full model: excess parameters also decay); shadow value of a unique token at \$$10^{-18}$ per FLOP; $^c\hat w/(1+m_N)$: the wedge measured from the capped choice relative to one plus the value of compactness there ($\hat w/w$ is 0.86 in 2029 on the rising path); $\sigma^*=0.60$, 0.74: same path and frontier as the reference; frontier rows: $5.06\times$.
\end{tablenotes}
\begin{tablenotes}[Source]
\citet{epochai2026data}; \citet{sevilla2024training}; \citet{hoffmann2022training}; \citet{besiroglu2024chinchilla}; \citet{grattafiori2024llama}; \citet{li2025predictableb}; \citet{gadre2024language}; \citet{bhagia2024establishing}; \citet{villalobos2022run}; \citet{muennighoff2023scaling}; author's calculations.
\end{tablenotes}
\end{table}


\subsection{Growth Models}\label{sec:econ-growth}

For models of AI and growth that map compute into capability, the relevant elasticity is $\gamma$, the compute
elasticity of reducible loss along the frontier; near the expansion path compute is a sufficient statistic for the
input bundle, and $\sigma^*$ enters only when the input mix is constrained, as behind a data wall. For calibrations at
$10^{25}$ to $10^{27}$ FLOP we recommend the reference value, $\gamma=0.165$ (standard error 0.006). Halving reducible
loss then takes 66 times the compute (95 percent confidence interval 50 to 90), or 2.6 years of frontier compute growth
at 5.06-fold a year and 2.9 years at 4.2-fold. Because $\gamma$ is cardinal in reducible loss, it needs the irreducible
loss $E$, and the Chinchilla design \citep{hoffmann2022training} alone has known loss units (it defines its loss as
per-token log loss on MassiveText, that is, nats per token) and identifies $E$ twice, from off-path runs and from its
IsoFLOP minima (Online Appendix~\ref{app:econ}); unlike $\sigma^*$, its local $\gamma$ does not drift with compute.
Other fits, of Chinchilla with $\kappa=1$ and of the Llama~3 and Farseer designs, give 0.14 to 0.18, or 49 to 161 times
the compute per halving (Online Appendix Table~\ref{tab:app-econ-gamma}). The recommendation extrapolates three to five decades beyond the design, and loss is not capability: at $10^{26}$ FLOP
reducible loss is about 4 percent of total loss in Chinchilla's units, so a calibration must say which loss it maps into
output. Integrated-assessment models such as GATE \citep{erdil2025gate} can use these numbers only on their input
side.
